% © 2026 Ing. David Jaroš — CC BY-NC-ND 4.0
%
% This work is licensed under a Creative Commons Attribution-NonCommercial-NoDerivatives
% 4.0 International License (CC BY-NC-ND 4.0).
%
% License History: Earlier drafts (up to v0.3) were released under CC BY 4.0.
% From v0.4 onward, all material is released under CC BY-NC-ND 4.0 to protect
% the integrity of the theoretical work during ongoing academic development.

\documentclass[11pt,a4paper]{article}
\usepackage{amsmath, amssymb, amsthm}
\usepackage{hyperref}
\usepackage{geometry}
\geometry{margin=2.5cm}

\newtheorem{theorem}{Theorem}
\newtheorem{proposition}{Proposition}
\newtheorem{definition}{Definition}
\newtheorem{remark}{Remark}
\newtheorem{conjecture}{Conjecture}

\title{Torus Geometry and Complex Time in UBT\\
\large Formal Relationship between $\tau$, Modular Structure, and the Upper Half-Plane}
\author{Ing.\ David Jaroš}
\date{2026}

\begin{document}
\maketitle

\begin{abstract}
The complex time $\tau = t + i\psi$ in the Unified Biquaternion Theory (UBT)
defines a natural torus geometry $\mathbb{C}/(\mathbb{Z} + \tau\mathbb{Z})$
when the imaginary direction $\psi$ is compact. This document formalises the
relationship between the UBT parameter $\tau$, the upper half-plane $\mathbb{H}$,
the modular group $\mathrm{SL}(2,\mathbb{Z})$, and the $q$-parameter used in
modular form expansions. The relationship between the compactification radius
$R_\psi$ and $\mathrm{Im}(\tau)$ is analysed. All claims are labelled with
their proof status.
\end{abstract}

% ======================================================================
\section{Complex Time: Definition and Compactification}
\label{sec:tau_def}
% ======================================================================

\begin{definition}[Complex Time]\label{def:complex_time}
The complex time coordinate in UBT is:
\begin{equation}
  \tau \;=\; t + i\psi, \qquad t \in \mathbb{R},\quad \psi \in \mathbb{R}.
\end{equation}
This is Axiom A2 of UBT. See \texttt{THEORY/canonical/canonical\_axioms.tex}, §A2.
\end{definition}

The imaginary component $\psi$ is compactified:
\begin{equation}
  \psi \;\sim\; \psi + 2\pi R_\psi,
\end{equation}
where $R_\psi > 0$ is the \emph{compactification radius}. This is a free
parameter of UBT at the current stage of development; see
\texttt{AUDITS/complex\_time\_modular\_audit.md} §3 for the status of
attempts to fix $R_\psi$ from first principles.

\begin{remark}[Status]
The compactification of $\psi$ follows from requiring a finite Kaluza--Klein
spectrum and from the periodicity of the Θ-field on the imaginary axis.
Status: \textbf{[L1]} — proved given the compactness assumption.
Source: \texttt{report/theta\_spectrum\_analysis.md} §2.1.
\end{remark}

% ======================================================================
\section{Torus Construction}
\label{sec:torus}
% ======================================================================

\begin{definition}[Torus from $\tau$]\label{def:torus}
Let $\tau_{\mathrm{mod}} \in \mathbb{H}$ be defined by:
\begin{equation}
  \tau_{\mathrm{mod}} \;=\; i\,\frac{R_\psi}{R_t},
\end{equation}
where $R_t > 0$ is a real-time scale (taken equal to $R_\psi$ in the
isotropic case $R_t = R_\psi$, giving $\tau_{\mathrm{mod}} = i$).
The UBT torus is the complex torus:
\begin{equation}
  T^2(\tau_{\mathrm{mod}}) \;=\; \mathbb{C}\big/\!\left(\mathbb{Z} + \tau_{\mathrm{mod}}\,\mathbb{Z}\right).
\end{equation}
\end{definition}

\begin{remark}[Status of $\tau_{\mathrm{mod}}$]
The identification of $\tau_{\mathrm{mod}} = iR_\psi/R_t$ is \textbf{[L1]}:
it follows from the lattice $\Lambda = R_t\,\mathbb{Z} \oplus R_\psi\,\mathbb{Z}$
generated by the two compact directions. The question of whether $R_t$ can be
fixed from UBT is \textbf{[O]} (open); $R_t$ enters only as an overall scale
if the theory is scale-invariant.
\end{remark}

% ======================================================================
\section{Upper Half-Plane and Modular Group}
\label{sec:upper_half_plane}
% ======================================================================

\begin{definition}[Upper Half-Plane]
\begin{equation}
  \mathbb{H} \;=\; \{\,\tau \in \mathbb{C} : \mathrm{Im}(\tau) > 0\,\}.
\end{equation}
\end{definition}

For UBT, $\tau_{\mathrm{mod}} = iR_\psi/R_t$ satisfies $\mathrm{Im}(\tau_{\mathrm{mod}}) = R_\psi/R_t > 0$,
so $\tau_{\mathrm{mod}} \in \mathbb{H}$ (\textbf{[L0]}).

\begin{definition}[Modular Group]
\begin{equation}
  \mathrm{SL}(2,\mathbb{Z}) \;=\; \left\{\, \begin{pmatrix} a & b \\ c & d \end{pmatrix} :
    a,b,c,d \in \mathbb{Z},\; ad - bc = 1 \right\}.
\end{equation}
This group acts on $\mathbb{H}$ by Möbius transformations:
\begin{equation}
  \tau \;\mapsto\; \frac{a\tau + b}{c\tau + d}.
\end{equation}
The generators are the T-transformation $\tau \mapsto \tau + 1$ and the
S-transformation $\tau \mapsto -1/\tau$.
\end{definition}

\begin{proposition}[T-symmetry of UBT]\label{prop:T_sym}
The UBT Θ-field is invariant (up to a phase) under $\tau \mapsto \tau + 1$.

\noindent\textbf{Status:} \textbf{[L1]}.
\noindent\textbf{Source:} \texttt{report/hecke\_modular\_structure.md} §2.1.
\end{proposition}

\begin{proposition}[S-symmetry breaking]\label{prop:S_break}
The S-transformation $\tau_{\mathrm{mod}} \mapsto -1/\tau_{\mathrm{mod}}$ maps
$iR_\psi/R_t \mapsto iR_t/R_\psi$, i.e.\ it exchanges the two radii.
This is \textbf{T-duality} on the $\psi$-circle. The Θ-field zero mode
(Jacobi theta constant $\vartheta_3$) transforms under S as:
\begin{equation}
  \vartheta_3(0|\tau) \;=\; (-i\tau)^{1/2}\,\vartheta_3\!\left(0\,\Big|\,-\frac{1}{\tau}\right),
\end{equation}
which is \emph{not} an invariance. The UBT imaginary axis $\tau \in i\mathbb{R}_{>0}$
is not preserved by the full modular group.

\noindent\textbf{Status:} \textbf{[L1]} for the T-duality statement;
\textbf{[L0]} for the theta transformation formula (classical result).
\noindent\textbf{Source:} \texttt{report/hecke\_modular\_structure.md} §2.1;
Jacobi 1829.
\end{proposition}

% ======================================================================
\section{The $q$-Parameter}
\label{sec:q_parameter}
% ======================================================================

\begin{definition}[$q$-parameter]
Given $\tau_{\mathrm{mod}} \in \mathbb{H}$, the $q$-parameter (nome) is:
\begin{equation}
  q \;=\; e^{i\pi\tau_{\mathrm{mod}}} \;=\; e^{-\pi R_\psi/R_t}.
\end{equation}
For the isotropic case $R_\psi = R_t$: $q = e^{-\pi} \approx 0.0432$.
\end{definition}

\begin{remark}[Role of $\mathrm{Im}(\tau)$]
The quantity $|q| = e^{-\pi\,\mathrm{Im}(\tau_{\mathrm{mod}})}$ controls
\emph{spectral suppression}: higher KK modes are suppressed as $q^{n^2}$
in the Jacobi theta expansion
\begin{equation}
  \vartheta_3(0|\tau_{\mathrm{mod}}) \;=\; \sum_{n \in \mathbb{Z}} q^{n^2}
    \;=\; 1 + 2q + 2q^4 + 2q^9 + \cdots
\end{equation}
For large $\mathrm{Im}(\tau)$ (i.e.\ large $R_\psi/R_t$), higher modes are
exponentially suppressed.
\end{remark}

% ======================================================================
\section{Role of $\mathrm{Im}(\tau)$ and the Compactification Scale}
\label{sec:im_tau}
% ======================================================================

\subsection{$\mathrm{Im}(\tau)$ as Physical Compactification Scale}

The relation $\mathrm{Im}(\tau_{\mathrm{mod}}) = R_\psi/R_t$ shows that
$\mathrm{Im}(\tau)$ directly encodes the ratio of compactification radii.
Large $\mathrm{Im}(\tau)$ corresponds to a large imaginary direction relative
to the real-time scale.

\subsection{Can Modular Minima Fix $\mathrm{Im}(\tau)$?}

The Dedekind eta function:
\begin{equation}
  \eta(\tau) \;=\; q^{1/24}\prod_{n=1}^{\infty}(1 - q^{2n}), \qquad q = e^{i\pi\tau},
\end{equation}
satisfies $|\eta(\tau)| \to 0$ as $\mathrm{Im}(\tau) \to \infty$ and has no
minimum on the imaginary axis $\tau = iy$, $y > 0$.

The effective potential $V_{\mathrm{eff}}$ studied in the $\alpha$ derivation
program does exhibit a minimum at $n^* = 137$ among primes, but this minimum
is in the discrete $n$-space, not in $\tau$-space. The \emph{modular}
interpretation of this minimum is:

\begin{equation}
  \tau^* \;\stackrel{?}{=}\; \frac{i}{n^*} \;=\; \frac{i}{137},
\end{equation}

giving $\mathrm{Im}(\tau^*) = 1/137 = \alpha_0$ (the fine-structure constant
at tree level). This identification is \textbf{[SE]}: it is consistent with
the V_eff calculation, but the \emph{B\_base} coefficient entering $V_{\rm eff}$
is not yet derived. See \texttt{research/theta\_alpha\_connection.md} §3.1.

\subsection{Reduction of Arbitrariness}

If $\mathrm{Im}(\tau)$ were fixed by a modular extremum, the single remaining
free parameter of the $\alpha$ program would be eliminated. The current status is:

\begin{itemize}
  \item V_eff minimum identifies $n^* = 137$: \textbf{[L1]} (proved given $B_{\rm base}$)
  \item Modular interpretation $\tau^* = i/n^*$: \textbf{[SE]}
  \item B_base derivation: \textbf{[O]} (open hard problem)
\end{itemize}

% ======================================================================
\section{Summary and Cross-References}
\label{sec:summary}
% ======================================================================

\begin{center}
\begin{tabular}{|l|l|l|}
\hline
\textbf{Object} & \textbf{UBT Status} & \textbf{Primary Source} \\
\hline
$\tau = t + i\psi$ & [L0] Axiom A2 & \texttt{canonical\_axioms.tex} \\
$\psi \sim \psi + 2\pi R_\psi$ & [L1] & \texttt{theta\_spectrum\_analysis.md} \\
$T^2(\tau) = \mathbb{C}/(\mathbb{Z}+\tau\mathbb{Z})$ & [L1] & This document \\
$\tau_{\rm mod} \in \mathbb{H}$ & [L0] & This document \\
T-symmetry $\tau \mapsto \tau+1$ & [L1] & \texttt{hecke\_modular\_structure.md} \\
S-symmetry breaking & [L1] & \texttt{hecke\_modular\_structure.md} \\
$q = e^{i\pi\tau}$ & [L0] & Classical; this document \\
$\vartheta_3(0|\tau)$ = Θ zero mode & [L1] & \texttt{hecke\_modular\_structure.md} \\
$\mathrm{Im}(\tau) \to \alpha_0$? & [SE] & \texttt{theta\_alpha\_connection.md} \\
B\_base derivation & [O] & \texttt{STATUS\_ALPHA.md} \\
\hline
\end{tabular}
\end{center}

\bigskip

\noindent\textbf{Related documents:}
\begin{itemize}
  \item \texttt{AUDITS/complex\_time\_modular\_audit.md} — full audit
  \item \texttt{research/theta\_alpha\_connection.md} — alpha candidate mechanisms
  \item \texttt{research/moduli\_space\_ads\_vs\_physical\_ds.tex} — AdS vs.\ dS distinction
  \item \texttt{research/mirror\_sector\_modular\_status.md} — mirror sector
  \item \texttt{report/hecke\_modular\_structure.md} — Hecke operator analysis
  \item \texttt{consolidation\_project/appendix\_ALPHA\_torus\_theta.tex} — eta/torus alpha
\end{itemize}

\end{document}
